% theme: area_bisecting_lines
\MajorQuestionLesson{二等分線の利用問題}
\SetRowSpace{4mm}

% 三角形：頂点を通る二等分線
\TypeQuestion{}{figure_application_triangle_median}
\FigProblemBlock{45mm}{fig/lesson_07_triangle_median}{\TextTall{図のように、放物線 $y=x^2$ 上に2点 $\mathrm{A},\mathrm{B}$ があり、それぞれの $x$ 座標は $-1,3$ である。また、点 $\mathrm{C}$ の座標は $(-1,-3)$ である。点 $\mathrm{A}$ を通り、$\Tri{ABC}$ の面積を2等分する直線を $\ell$ とする。次の問いに答えなさい。}}{
\QQ{点 $\mathrm{A},\mathrm{B}$ の座標をそれぞれ求めなさい。}{}
\QQ{辺 $\mathrm{BC}$ の中点 $\mathrm{M}$ の座標を求めなさい。}{}
\QQ{直線 $\ell$ の式を求めなさい。}{}
\QQ{$\Tri{ABC}$ の面積と、直線 $\ell$ によって分けられた2つの三角形の面積をそれぞれ求めなさい。}{}
}

% 三角形：頂点を通らない二等分線
\TypeQuestion{}{figure_application_triangle_bisector_without_vertex}
\FigProblemBlock{45mm}{fig/lesson_07_triangle_no_vertex}{\TextTall{図のように、放物線 $y=x^2$ 上に2点 $\mathrm{A},\mathrm{B}$ があり、それぞれの $x$ 座標は $-3,3$ である。また、点 $\mathrm{C}$ の座標は $(1,-3)$ である。辺 $\mathrm{AB}$ 上に $x$ 座標が $1$ である点 $\mathrm{D}$ をとる。辺 $\mathrm{AC}$ 上に点 $\mathrm{E}$ をとり、線分 $\mathrm{DE}$ が $\Tri{ABC}$ の面積を2等分するとき、次の問いに答えなさい。}}{
\QQ{点 $\mathrm{A},\mathrm{B},\mathrm{D}$ の座標をそれぞれ求めなさい。}{}
\QQ{直線 $\mathrm{AC}$ の式を求めなさい。}{}
\QQ{$\Tri{ABC}$ の面積を求めなさい。}{}
\QQ{点 $\mathrm{E}$ の座標を求めなさい。}{}
\QQ{直線 $\mathrm{DE}$ の式を求めなさい。}{}
}

\LessonPageBreak

% 平行四辺形：下に凸
\TypeQuestion{}{figure_application_parallelogram_bisection}
\FigProblemBlock{45mm}{fig/lesson_07_parallelogram_down}{\TextTall{図のように、放物線 $y=x^2$ 上に2点 $\mathrm{A},\mathrm{B}$ があり、それぞれの $x$ 座標は $-2,2$ である。点 $\mathrm{D}$ は原点であり、四角形 $\mathrm{ABCD}$ は平行四辺形である。対角線 $\mathrm{AC},\mathrm{BD}$ の交点を $\mathrm{P}$ とする。また、点 $\mathrm{P}$ を通る直線 $\ell:y=-x+k$ が平行四辺形 $\mathrm{ABCD}$ の面積を2等分している。次の問いに答えなさい。}}{
\QQ{点 $\mathrm{A},\mathrm{B},\mathrm{C}$ の座標をそれぞれ求めなさい。}{}
\QQ{点 $\mathrm{P}$ の座標を求めなさい。}{}
\QQ{定数 $k$ の値を求めなさい。}{}
\QQ{直線 $\ell$ と辺 $\mathrm{AB},\mathrm{CD}$ との交点をそれぞれ $\mathrm{E},\mathrm{F}$ とする。点 $\mathrm{E},\mathrm{F}$ の座標を求めなさい。}{}
\QQ{直線 $\ell$ によって分けられた2つの部分の面積を、それぞれ求めなさい。}{}
}

% 平行四辺形：上に凸・傾きが変わる直線
\TypeQuestion{}{figure_application_parallelogram_bisection_upper}
\FigProblemBlock{45mm}{fig/lesson_07_parallelogram_up}{\TextTall{図のように、放物線 $y=-x^2$ 上に2点 $\mathrm{A},\mathrm{B}$ があり、それぞれの $x$ 座標は $-2,2$ である。点 $\mathrm{D}$ は原点であり、四角形 $\mathrm{ABCD}$ は平行四辺形である。対角線 $\mathrm{AC},\mathrm{BD}$ の交点を $\mathrm{P}$ とする。また、点 $\mathrm{P}$ を通る直線 $m:y=ax-3$ が平行四辺形 $\mathrm{ABCD}$ の面積を2等分している。次の問いに答えなさい。}}{
\QQ{点 $\mathrm{A},\mathrm{B},\mathrm{C}$ の座標をそれぞれ求めなさい。}{}
\QQ{点 $\mathrm{P}$ の座標を求めなさい。}{}
\QQ{定数 $a$ の値を求めなさい。}{}
\QQ{直線 $m$ と辺 $\mathrm{AB},\mathrm{CD}$ との交点をそれぞれ $\mathrm{E},\mathrm{F}$ とする。点 $\mathrm{E},\mathrm{F}$ の座標を求めなさい。}{}
\QQ{直線 $m$ によって分けられた2つの部分の面積を、それぞれ求めなさい。}{}
}
